\documentclass[12pt, fleqn]{article}

\usepackage{../../../template/template}
\usepackage{../../../template/fortickets}
%\usepackage{../../../template/KillContents}

\begin{document}
  \begin{figure}[H]
      \includegraphics[width=10cm]{../../../template/spsu.jpg}
      \centering
  \end{figure}
  \vspace*{\fill}
  \begin{center}
    \huge Билеты по дискре, 1-2 сем

    \Large (преподаватель Григорьева Н. С.)\footnote{Также были использованы учебник Романовский И.В. "Дискретный анализ"{} и интернет}

    \large Записали Костин П.А. и Щукин И.В.\footnote{Данный документ неидеальный, прошу сообщать о найденных недочетах \href{https://vk.com/drab_existence_a}{вконтакте} (можно присылать скрины неточностей с указанием билетов)}
  \end{center}
  \vspace*{\fill}
  \vspace*{\fill}
  
  \newpage
  \tableofcontents
  \newpage

  %\section{Некоторые определения из теории множеств. Прямое произведение, разбиение множеств. Мощность объединения}
  \subfile{1.tex}

  %\section{Вектора из нулей и единиц}
  \subfile{2.tex}

  %\section{Алгоритм перебора 0-1 векторов. Коды Грея}
  \subfile{3.tex}

  %\section{Перебор элементов прямого произведения множеств}
  \subfile{4.tex}

  %\section{Размещения, сочетания, перестановки без повторений}
  \subfile{5.tex}

  %\section{Размещения, сочетания, перестановки с повторениями}
  \subfile{6.tex}

  %\section{Два алгоритма перебора перестановок. Нумерация перестановок}
  \subfile{7.tex}

  %\section{Задача о минимуме скалярного произведения}
  \subfile{8.tex}

  %\section{Числа Фибоначчи. Теорема о представлении}
  \subfile{9.tex}

  %\section{Перебор сочетаний. Нумерация сочетаний}
  \subfile{10.tex}

  %\section{Бином Ньютона и его комбинаторное использование}
  \subfile{11.tex}

  %\section{Свойства биномиальных коэффициентов}
  \subfile{12.tex}

  %\section{Основные определения теории вероятностей}
  \subfile{13.tex}

  %\section{Условные вероятности и формула Байеса}
  \subfile{14.tex}

  %\section{Математическое ожидание и дисперсия случайной величины}
  \subfile{15.tex}

  %\section{Схема Бернулли}
  \subfile{16.tex}

  %\section{Случайные числа. Схема Уолкера}
  \subfile{17.tex}

  %\section{Двоичный поиск и неравенство Крафта}
  \subfile{18.tex}

  %\section{Энтропия. 2 леммы}
  \subfile{19.tex}

  %\section{Теорема об энтропии}
  \subfile{20.tex}

  %\section{Операции над строками переменной длины}
  \subfile{21.tex}

  %\section{Поиск образца в строке (Карпа-Рабина, Бойера-Мура)}
  \subfile{22.tex}

  %\section{Суффиксное дерево}
  \subfile{23.tex}

  %\section{Задача о максимальном совпадении двух строк}
  \subfile{24.tex}

  %\section{Код Шеннона-Фано. Алгоритм Хаффмена. 3 леммы}
  \subfile{25.tex}

  %\section{Сжатие информации по методу Зива-Лемпеля}
  \subfile{26.tex}

  %\section{Метод Барроуза-Уилера}
  \subfile{27.tex}

  \section{Избыточное кодирование. Коды Хэмминга}
  \newpage

  \section{Шифрование с открытым ключом}
  \newpage

  \section{Сортировки (5 методов)}
  \newpage

  \section{Информационный поиск и организация информации}
  \newpage

  \section{Хеширование}
  \newpage

  %\section{АВЛ-деревья}
  \subfile{33.tex}

  \section{B-деревья}
  \newpage

  %\section{Биноминальные кучи}
  \subfile{35.tex}

  \section{Основные определения теории графов}
  \newpage

  \section{Построение транзитивного замыкания}
  \newpage

  \section{Обходы графа в ширину и глубину. Топологическая сортировка}
  \newpage

  \section{Связность. Компоненты связности и сильной связности}
  \newpage

  \section{Алгоритм поиска контура и построение диаграммы порядка}
  \newpage

  \section{Теорема о связном подграфе}

  \subfile{41.tex}
  \newpage

  \section{Деревья. Теорема о шести эквивалентных определениях дерева}
  \subfile{42.tex}
  \newpage

  \section{Задача о кратчайшем остовном дереве. Алгоритм Прима}
  \newpage

  \section{Алгоритм Краскала}
  \newpage

  \section{Задача о кратчайшем пути. Алгоритм Дейкстры}
  \newpage

  \section{Алгоритм Левита}
  \subfile{46.tex}
  \newpage

  \section{Задача о кратчайшем дереве путей}
  \subfile{47.tex}
  \newpage

  \section{Сетевой график и критические пути. Нахождение резервов работ}
  \newpage

  \section{Задача о максимальном паросочетании в графе. Алгоритм построения}
  \subfile{49.tex}
  \newpage

  \section{Теорема Кенига}
  \subfile{50.tex}
  \newpage

  \section{Алгоритм построения контролирующего множества}
  \newpage

  \subfile{52.tex}
  \newpage

  \section{Задача коммивояжера. Метод ветвей и границ}
  \newpage

  \section{Метод динамического программирования. Задача линейного раскроя}


  \section{Приближенные методы решения дискретных задач. Жадные алгоритмы}


  \section{Алгоритмы с гарантированной оценкой точности. Алгоритм Эйлера}


  \section{Жадные алгоритмы. Задача о системе различных представителей}


  \section{Приближенные методы решения дискретных задач}


  \section{Конечные автоматы}


  \section{Числа Фибоначчи. Производящие функции}


  \section{Числа Каталана}


  \section{?Алгоритм Кристофидеса (возможно будет)}

\end{document}
